\documentclass[10pt,aspectratio=169,serif]{beamer}

\usepackage[fontset=none, scheme=plain]{ctex}
\IfFontExistsTF{Songti SC}{%
	\setCJKmainfont{Songti SC}[AutoFakeSlant=0.2]%
	\setCJKsansfont{Songti SC}[AutoFakeSlant=0.2]%
}{%
	\setCJKmainfont{FandolSong-Regular}[Extension=.otf,BoldFont=FandolSong-Bold,AutoFakeSlant=0.2]%
	\setCJKsansfont{FandolSong-Regular}[Extension=.otf,BoldFont=FandolSong-Bold]%
}
\usepackage{amsmath,amssymb,graphicx,tikz,booktabs,multicol,calligra}
\usepackage{pgfplots}
\pgfplotsset{compat=1.18}
\usepackage{tcolorbox}
\tcbuselibrary{skins}
\usepackage{XyygU}
\usepackage{caption}
\captionsetup[figure]{name=图,labelsep=period}

\usetikzlibrary{arrows.meta,positioning,calc}

\setlength{\abovedisplayskip}{0pt plus 2pt minus 2pt}
\setlength{\belowdisplayskip}{0pt plus 2pt minus 2pt}

\title[习题4.4]{习题4.4 对数函数}
\subtitle{复习巩固 \& 综合运用 \& 拓广探索}
\author{}
\institute{}
\date{2026-10-16}

\begin{document}

\begin{frame}
  \titlepage
\end{frame}

% ========== 复习巩固 第1题 ==========
\begin{frame}[t]{第1题：求函数定义域}
\textbf{1.} 求下列函数的定义域：

\vspace{1em}
(1) $y = \sqrt[3]{\log_2 x}$；

\vspace{1em}
(2) $y = \sqrt{\log_{0.5}(4x-3)}$。
\end{frame}

\begin{frame}[t]{第1题 解答}
\begin{tcolorbox}[colback=orange!10,colframe=orange,title=\textbf{解}]
(1) 三次根号对任意实数有意义，只需 $\log_2 x$ 有意义，即 $x>0$。

\hspace{2em}定义域：$\alert{(0, +\infty)}$。

(2) 需同时满足：
\[
\begin{cases}
4x-3 > 0,\\
\log_{0.5}(4x-3) \geq 0.
\end{cases}
\]
由底数 $0.5<1$，$\log_{0.5}(4x-3)\geq 0 = \log_{0.5} 1$，得 $4x-3 \leq 1$，即 $x\leq 1$。

\hspace{2em}又 $x>\dfrac{3}{4}$，故定义域：$\alert{\left(\dfrac{3}{4},\, 1\right]}$。
\end{tcolorbox}
\end{frame}

% ========== 第2题 ==========
\begin{frame}[t]{第2题：比较正数 $m, n$ 的大小}
\textbf{2.} 比较满足下列条件的两个正数 $m, n$ 的大小：

\vspace{0.5em}
\begin{columns}[t]
\begin{column}{0.5\textwidth}
\begin{enumerate}
  \item $\log_3 m < \log_3 n$；
  \item $\log_{0.3} m < \log_{0.3} n$；
\end{enumerate}
\end{column}
\begin{column}{0.5\textwidth}
\begin{enumerate}\setcounter{enumi}{2}
  \item $\log_a m < \log_a n$（$0<a<1$）；
  \item $\log_a m > \log_a n$（$a>1$）。
\end{enumerate}
\end{column}
\end{columns}
\end{frame}

\begin{frame}[t]{第2题 解答}
\begin{tcolorbox}[colback=orange!10,colframe=orange,title=\textbf{解}]
(1) 底数 $3>1$，$y=\log_3 x$ 是增函数，由 $\log_3 m<\log_3 n$ 得 $\alert{m<n}$。

(2) 底数 $0.3<1$，$y=\log_{0.3} x$ 是减函数，由 $\log_{0.3} m<\log_{0.3} n$ 得 $\alert{m>n}$。

(3) $0<a<1$，$y=\log_a x$ 是减函数，由 $\log_a m<\log_a n$ 得 $\alert{m>n}$。

(4) $a>1$，$y=\log_a x$ 是增函数，由 $\log_a m>\log_a n$ 得 $\alert{m>n}$。
\end{tcolorbox}
\end{frame}

% ========== 第3题 ==========
\begin{frame}[t]{第3题：火箭速度}
\textbf{3.} 在不考虑空气阻力的条件下，火箭的最大速度 $v$（单位：km/s）和燃料质量 $M$（单位：kg）、火箭（除燃料外）的质量 $m$（单位：kg）的函数关系是
\[
v = 2000\ln\left(1+\frac{M}{m}\right).
\]
当燃料质量是火箭质量的百分之几时，火箭的最大速度可达到 $12$ km/s？
\end{frame}

\begin{frame}[t]{第3题 解答}
\begin{tcolorbox}[colback=orange!10,colframe=orange,title=\textbf{解}]
令 $v=12$：
\[
2000\ln\left(1+\frac{M}{m}\right) = 12 \;\Longrightarrow\; \ln\left(1+\frac{M}{m}\right) = 0.006.
\]
所以：
\[
1+\frac{M}{m} = e^{0.006} \approx 1.00602.
\]
\[
\frac{M}{m} \approx 0.00602 \approx 0.6\%.
\]
即当燃料质量约是火箭质量的 $\alert{0.6\%}$ 时，火箭最大速度可达 $12$ km/s。
\end{tcolorbox}
\end{frame}

% ========== 第4题 ==========
\begin{frame}[t]{第4题：对数函数图象识别}
\textbf{4.} 函数 $y=\log_2 x$，$y=\log_5 x$，$y=\lg x$ 的图象如图所示（三条曲线都过 $(1,0)$）。

\vspace{0.3em}
(1) 试说明哪个函数对应于哪个图象，并解释为什么；

(2) 在同一直角坐标系中画出 $y=\log_{\frac{1}{2}} x$，$y=\log_{\frac{1}{5}} x$，$y=\lg_{\frac{1}{10}} x$ 的图象；

(3) 从(2)的图中你发现了什么？
\end{frame}

\begin{frame}[t]{第4题 解答}
\begin{tcolorbox}[colback=orange!10,colframe=orange,title=\textbf{解}]
(1) 当 $x>1$ 时，底数越小，对数值越大（$\log_a x = \dfrac{\lg x}{\lg a}$，$0<\lg a<1$，$a$ 越小则 $\lg a$ 越小）。

\hspace{2em}所以 $\log_2 x > \lg x > \log_5 x$（$x>1$ 时）。

\hspace{2em}图中③最高，②次之，①最低，故：$\alert{③=\log_2 x,\;②=\lg x,\;①=\log_5 x}$。

(2) 因为 $\log_{1/a} x = -\log_a x$，所以它们的图象分别与(1)中对应图象关于 $x$ 轴对称。

(3) 互为倒数的两个底数（如 $2$ 与 $\frac{1}{2}$），其对数函数图象关于 $x$ 轴对称。
\end{tcolorbox}
\end{frame}

\begin{frame}[t]{第4题 图象}
\begin{center}
\begin{tikzpicture}[scale=0.75]
\begin{axis}[
    width=10cm, height=6cm,
    axis lines=middle,
    xlabel=$x$, ylabel=$y$,
    xmin=0, xmax=10, ymin=-3, ymax=3,
    xtick={1,5,10}, ytick={-2,-1,1,2},
    grid=major, grid style={gray!20},
]
\addplot[thick,red,domain=0.05:10,samples=100] {ln(x)/ln(2)};
\node[red] at (axis cs:6,2) {③ $\log_2 x$};
\addplot[thick,blue,domain=0.05:10,samples=100] {ln(x)/ln(10)};
\node[blue] at (axis cs:6,1) {② $\lg x$};
\addplot[thick,green!60!black,domain=0.05:10,samples=100] {ln(x)/ln(5)};
\node[green!60!black] at (axis cs:6,0.5) {① $\log_5 x$};
\end{axis}
\end{tikzpicture}
\end{center}
\end{frame}

% ========== 第5题 ==========
\begin{frame}[t]{第5题：鲑鱼游速}
\textbf{5.} 大西洋鲑鱼的游速 $v$（单位：m/s）可以表示为 $v = \dfrac{1}{2}\log_3\dfrac{O}{100}$，其中 $O$ 表示鱼的耗氧量的单位数。

\vspace{0.5em}
(1) 当一条鱼的耗氧量是 $2700$ 个单位时，它的游速是多少？

(2) 计算一条鱼静止时耗氧量的单位数。
\end{frame}

\begin{frame}[t]{第5题 解答}
\begin{tcolorbox}[colback=orange!10,colframe=orange,title=\textbf{解}]
(1) 当 $O=2700$ 时：
\[
v = \frac{1}{2}\log_3\frac{2700}{100} = \frac{1}{2}\log_3 27 = \frac{1}{2}\times 3 = \alert{1.5\ \text{m/s}}.
\]

(2) 静止时 $v=0$：
\[
\frac{1}{2}\log_3\frac{O}{100} = 0 \;\Longrightarrow\; \log_3\frac{O}{100} = 0 \;\Longrightarrow\; \frac{O}{100} = 1.
\]
所以 $\alert{O = 100}$ 个单位。
\end{tcolorbox}
\end{frame}

% ========== 第6题 ==========
\begin{frame}[t]{第6题：药物含量图象选择}
\textbf{6.} 在 $2$\,h 内将某种药物注射进患者血液。注射期间，血液中药物含量呈线性增加；停止注射后，药物含量呈指数衰减。能反映血液中药物含量 $Q$ 随时间 $t$ 变化的图象是（\quad）。

\vspace{0.5em}
\small（A）：$0\sim 1$\,h 线性上升，$1$\,h 后指数衰减；
（B）：$0\sim 2$\,h 线性上升，$2$\,h 后指数衰减；

\small（C）：从 $t=0$ 开始线性下降；
（D）：$0\sim 2$\,h 线性上升，$2$\,h 后变为负。
\end{frame}

\begin{frame}[t]{第6题 解答}
\begin{tcolorbox}[colback=orange!10,colframe=orange,title=\textbf{解}]
注射时间为 $2$\,h，此期间药物含量\alert{线性增加}，在 $t=2$\,h 达到峰值；
停止注射后药物\alert{指数衰减}（逐渐趋近于0，不会变为负值）。

\vspace{0.5em}
(A) 峰值在 $t=1$\,h，时间不对；
(C) 从 $t=0$ 就开始下降，与"注射期间线性增加"矛盾；
(D) 衰减后变为负值，药物含量不可能为负。

\vspace{0.3em}
\alert{答案：(B)}。
\end{tcolorbox}
\end{frame}

% ========== 第7题 ==========
\begin{frame}[t]{第7题：判断反函数}
\textbf{7.} 判断下列各对函数是否互为反函数。若是，则求出它们的定义域和值域：

\vspace{0.8em}
(1) $y = \ln x$，$y = e^x$；

\vspace{0.8em}
(2) $y = -\log_a x$，$y = \left(\dfrac{1}{a}\right)^x$。
\end{frame}

\begin{frame}[t]{第7题 解答}
\begin{tcolorbox}[colback=orange!10,colframe=orange,title=\textbf{解}]
(1) \alert{互为反函数}。

\hspace{2em}$y=\ln x$：定义域 $(0,+\infty)$，值域 $\mathbf{R}$；

\hspace{2em}$y=e^x$：定义域 $\mathbf{R}$，值域 $(0,+\infty)$。

(2) $y=-\log_a x = \log_a x^{-1} = \log_a\dfrac{1}{x} = \log_{1/a} x$，它是 $y=\left(\dfrac{1}{a}\right)^x$ 的反函数。

\hspace{2em}\alert{互为反函数}。

\hspace{2em}$y=-\log_a x$：定义域 $(0,+\infty)$，值域 $\mathbf{R}$；

\hspace{2em}$y=\left(\dfrac{1}{a}\right)^x$：定义域 $\mathbf{R}$，值域 $(0,+\infty)$。
\end{tcolorbox}
\end{frame}

% ========== 第8题 ==========
\begin{frame}[t]{第8题：摄氏与华氏温度}
\textbf{8.} 设 $y=f(x)$ 表示摄氏温度为 $x$ 时，华氏温度为 $y$。

\vspace{0.5em}
(1) 如果函数 $y=f(x)$ 的反函数是 $y=g(x)$，那么 $y=g(x)$ 表示什么？

(2) 如果 $f(30)=86$，那么求 $g(86)$，并说明其实际意义。
\end{frame}

\begin{frame}[t]{第8题 解答}
\begin{tcolorbox}[colback=orange!10,colframe=orange,title=\textbf{解}]
(1) $y=g(x)$ 表示：当华氏温度为 $x$ 时，对应的摄氏温度为 $y$。即 $g$ 是"由华氏温度求摄氏温度"的函数。

(2) 由反函数定义：若 $f(30)=86$，则 $g(86)=30$。

\hspace{2em}实际意义：华氏温度 $86^\circ\text{F}$ 对应摄氏温度 $30^\circ\text{C}$。
\end{tcolorbox}
\end{frame}

% ========== 第9题 ==========
\begin{frame}[t]{第9题：患病率降低}
\textbf{9.} 某地某种疾病的患病率正以每年 $15\%$ 的比例降低。要将当前的患病率降低一半，需要多少年？
\end{frame}

\begin{frame}[t]{第9题 解答}
\begin{tcolorbox}[colback=orange!10,colframe=orange,title=\textbf{解}]
设当前患病率为 $a$，经过 $x$ 年后患病率为原来的一半：
\[
a(1-15\%)^x = \frac{a}{2}.
\]
即 $0.85^x = 0.5$，两边取对数：
\[
x = \log_{0.85} 0.5 = \frac{\ln 0.5}{\ln 0.85} \approx \frac{-0.6931}{-0.1625} \approx 4.3.
\]
所以大约需要 $\alert{5}$ 年（向上取整）患病率降低一半。
\end{tcolorbox}
\end{frame}

% ========== 第10题 ==========
\begin{frame}[t]{第10题：声强级}
\textbf{10.} 声强级 $L_I$（单位：dB）由公式 $L_I = 10\lg\!\left(\dfrac{I}{10^{-12}}\right)$ 给出，其中 $I$ 为声强（单位：W/m$^2$）。

\vspace{0.5em}
(1) 一般正常人听觉能忍受的最高声强为 $1$\,W/m$^2$，能听到的最低声强为 $10^{-12}$\,W/m$^2$。求人听觉的声强级范围。

(2) 平时常人交谈时的声强约为 $10^{-6}$\,W/m$^2$，求其声强级。
\end{frame}

\begin{frame}[t]{第10题 解答}
\begin{tcolorbox}[colback=orange!10,colframe=orange,title=\textbf{解}]
(1) 当 $I=1$\,W/m$^2$ 时：
\[
L_I = 10\lg\frac{1}{10^{-12}} = 10\lg 10^{12} = 120\ \text{dB}.
\]
当 $I=10^{-12}$\,W/m$^2$ 时：
\[
L_I = 10\lg\frac{10^{-12}}{10^{-12}} = 10\lg 1 = 0\ \text{dB}.
\]
所以人听觉的声强级范围为 $\alert{0\sim 120\ \text{dB}}$。

(2) 当 $I=10^{-6}$\,W/m$^2$ 时：
\[
L_I = 10\lg\frac{10^{-6}}{10^{-12}} = 10\lg 10^6 = \alert{60\ \text{dB}}.
\]
\end{tcolorbox}
\end{frame}

% ========== 第11题 ==========
\begin{frame}[t]{第11题：房价增长方式}
\textbf{11.} 一套住房房价从2002年的 $20$ 万元涨到2012年的 $40$ 万元。$P_1$ 按直线上升，$P_2$ 按指数增长，$t$ 是2002年以来经过的年数。

\vspace{0.3em}
(1) 求 $P_1=f(t)$ 的解析式；

(2) 求 $P_2=g(t)$ 的解析式；

(3) 比较两种价格增长方式的差异。
\end{frame}

\begin{frame}[t]{第11题 解答}
\begin{tcolorbox}[colback=orange!10,colframe=orange,title=\textbf{解}]
(1) $P_1$ 是直线：过 $(0, 20)$ 和 $(10, 40)$，斜率 $k=\dfrac{40-20}{10}=2$。
\[
\alert{P_1 = f(t) = 20 + 2t}.
\]

(2) $P_2$ 是指数函数：过 $(0, 20)$ 和 $(10, 40)$，设 $P_2=20\cdot b^t$，则 $20\cdot b^{10}=40$，$b^{10}=2$，$b=2^{1/10}$。
\[
\alert{P_2 = g(t) = 20\cdot 2^{t/10}}.
\]

(3) 前 $10$ 年两者都从 $20$ 万涨到 $40$ 万；但长期来看，指数增长后期远远超过直线增长。如 $t=20$ 时，$P_1=60$ 万，$P_2=80$ 万。
\end{tcolorbox}
\end{frame}

% ========== 第12题 ==========
\begin{frame}[t]{第12题：求 $a$ 的取值范围}
\textbf{12.} 已知 $\log_a\dfrac{1}{2} < 1$，$\left(\dfrac{1}{2}\right)^a < 1$，$a^{\frac{1}{2}} < 1$，求实数 $a$ 的取值范围。
\end{frame}

\begin{frame}[t]{第12题 解答}
\begin{tcolorbox}[colback=orange!10,colframe=orange,title=\textbf{解}]
由 $\left(\dfrac{1}{2}\right)^a < 1 = \left(\dfrac{1}{2}\right)^0$，底数 $\frac{1}{2}<1$，得 $\alert{a > 0}$。

由 $a^{1/2} = \sqrt{a} < 1$，得 $\alert{0 < a < 1}$。

此时底数 $0<a<1$，$y=\log_a x$ 是减函数。由 $\log_a\dfrac{1}{2} < 1 = \log_a a$，得：
\[
\frac{1}{2} > a \quad(\text{减函数，不等号反向}).
\]

\vspace{0.3em}
综上：$\alert{0 < a < \dfrac{1}{2}}$。
\end{tcolorbox}
\end{frame}

% ========== 第13题 ==========
\begin{frame}[t]{第13题：比较三个值的大小}
\textbf{13.} 比较下列各题中三个值的大小：

\vspace{0.8em}
(1) $\log_{0.2} 6$，$\log_{0.3} 6$，$\log_{0.4} 6$；

\vspace{0.8em}
(2) $\log_2 3$，$\log_3 4$，$\log_4 5$。
\end{frame}

\begin{frame}[t]{第13题 解答}
\begin{tcolorbox}[colback=orange!10,colframe=orange,title=\textbf{解}]
(1) $\log_b 6 = \dfrac{1}{\log_6 b}$。因为 $0.2<0.3<0.4<1$，且 $\log_6 x$ 是增函数，所以 $\log_6 0.2 < \log_6 0.3 < \log_6 0.4 < 0$。

取倒数（负数取倒数，不等号反向）：
\[
\frac{1}{\log_6 0.2} > \frac{1}{\log_6 0.3} > \frac{1}{\log_6 0.4}.
\]
即 $\alert{\log_{0.2} 6 > \log_{0.3} 6 > \log_{0.4} 6}$。

(2) $\log_2 3 \approx 1.585$，$\log_3 4 \approx 1.262$，$\log_4 5 \approx 1.161$。

\vspace{0.3em}
所以 $\alert{\log_2 3 > \log_3 4 > \log_4 5}$。
\end{tcolorbox}
\end{frame}

\begin{frame}[c]
  \begin{center}
    {\Huge\calligra Thanks!}
  \end{center}
\end{frame}

\end{document}
